\section{Pengambilan sampel dari populasi kecil}
\label{smallPop}

\noindent%
Ketika kita mengambil pengamatan dari suatu populasi sebagai sampel,
biasanya kita hanya mengambil sebagian kecil dari
individu atau kasus yang mungkin.
Namun, terkadang ukuran sampel kita cukup besar
atau populasinya cukup kecil sehingga kita
mengambil lebih dari 10\% populasi\footnote{Pedoman 10\%
  merupakan batas berdasarkan kaidah praktis untuk menentukan kapan
  pertimbangan ini menjadi lebih penting.}
\emph{tanpa
pengembalian} (artinya kita tidak mungkin
mengambil kasus yang sama dua kali).
Pengambilan bagian populasi yang cukup besar seperti ini
dapat memengaruhi cara kita menganalisis
sampel tersebut.

\begin{examplewrap}
\begin{nexample}{Dosen terkadang memilih seorang mahasiswa secara acak untuk menjawab pertanyaan. Jika setiap mahasiswa memiliki peluang yang sama untuk dipilih dan ada 15 orang di kelas Anda, berapakah peluang dosen memilih Anda untuk pertanyaan berikutnya?}
Jika ada 15 orang yang dapat ditanya dan semuanya hadir di kelas, peluangnya adalah $1/15$, atau sekitar $0.067$.
\end{nexample}
\end{examplewrap}

\begin{examplewrap}
\begin{nexample}{Jika dosen mengajukan 3 pertanyaan, berapakah peluang Anda tidak terpilih? Anggap dosen tidak memilih orang yang sama dua kali dalam satu perkuliahan.}\label{3woRep}
Untuk pertanyaan pertama, dosen memilih orang lain dengan peluang $14/15$. Ketika mengajukan pertanyaan kedua, hanya ada 14 orang yang belum ditanya. Jadi, jika Anda tidak terpilih untuk pertanyaan pertama, peluang Anda kembali tidak terpilih adalah $13/14$. Demikian pula, peluang Anda kembali tidak terpilih untuk pertanyaan ketiga adalah $12/13$, dan peluang tidak terpilih untuk ketiga pertanyaan tersebut adalah
\begin{align*}
&P(\text{tidak terpilih dalam 3 pertanyaan}) \\
&\quad = P(\text{\var{Q1}} = \text{\resp{tidak\us{}terpilih}, }\text{\var{Q2}} = \text{\resp{tidak\us{}terpilih}, }\text{\var{Q3}} = \text{\resp{tidak\us{}terpilih}.}) \\
&\quad = \frac{14}{15}\times\frac{13}{14}\times\frac{12}{13} = \frac{12}{15} = 0.80
\end{align*}
\end{nexample}
\end{examplewrap}

\begin{exercisewrap}
\begin{nexercise}
Aturan apa yang memungkinkan kita mengalikan peluang-peluang dalam Contoh~\ref{3woRep}?\footnotemark
\end{nexercise}
\end{exercisewrap}
\footnotetext{Ketiga peluang yang kita hitung sebenarnya terdiri atas satu peluang marginal, $P($\var{Q1}$ = $\resp{tidak\us{}terpilih}$)$, dan dua peluang bersyarat:
\begin{align*}
&P(\text{\var{Q2}} =  \text{\resp{tidak\us{}terpilih} }|
    \text{ \var{Q1}} = \text{\resp{tidak\us{}terpilih}}) \\
&P(\text{\var{Q3}} =  \text{\resp{tidak\us{}terpilih} }|
    \text{ \var{Q1}} = \text{\resp{tidak\us{}terpilih}, }
    \text{\var{Q2}} = \text{\resp{tidak\us{}terpilih}})
\end{align*}
Dengan Aturan Perkalian Umum, hasil kali ketiga peluang ini adalah peluang tidak terpilih dalam 3 pertanyaan.}

\D{\newpage}

\begin{examplewrap}
\begin{nexample}{Misalkan dosen memilih secara acak tanpa memperhatikan siapa yang sudah dipilih, yaitu mahasiswa dapat dipilih lebih dari sekali. Berapakah peluang Anda tidak terpilih untuk ketiga pertanyaan tersebut?}\label{3wRep}
Pemilihan-pemilihan tersebut saling independen, dan peluang tidak terpilih untuk satu pertanyaan adalah $14/15$. Jadi, kita dapat menggunakan Aturan Perkalian untuk proses yang saling independen.
\begin{align*}
&P(\text{tidak terpilih dalam 3 pertanyaan}) \\
&\quad = P(\text{\var{Q1}} = \text{\resp{tidak\us{}terpilih}, }\text{\var{Q2}} = \text{\resp{tidak\us{}terpilih}, }\text{\var{Q3}} = \text{\resp{tidak\us{}terpilih}.}) \\
&\quad = \frac{14}{15}\times\frac{14}{15}\times\frac{14}{15} = 0.813
\end{align*}
Peluang Anda untuk tidak terpilih sedikit lebih besar dibandingkan ketika dosen memilih orang yang berbeda untuk setiap pertanyaan. Namun, sekarang Anda mungkin dipilih lebih dari sekali.
\end{nexample}
\end{examplewrap}

\begin{exercisewrap}
\begin{nexercise}
Berdasarkan skenario pada Contoh~\ref{3wRep}, berapakah peluang Anda terpilih untuk menjawab ketiga pertanyaan?\footnotemark
\end{nexercise}
\end{exercisewrap}
\footnotetext{$P($terpilih untuk menjawab ketiga pertanyaan$) = \left(\frac{1}{15}\right)^3 = 0.00030$.}

Jika kita mengambil sampel dari populasi kecil \term{tanpa pengembalian}, pengamatan-pengamatan kita tidak lagi saling independen. Dalam Contoh~\ref{3woRep}, peluang tidak terpilih untuk pertanyaan kedua dikondisikan pada peristiwa bahwa Anda tidak terpilih untuk pertanyaan pertama. Dalam Contoh~\ref{3wRep}, dosen mengambil sampel mahasiswa \term{dengan pengembalian}: dosen tersebut berulang kali mengambil sampel dari seluruh kelas tanpa memedulikan siapa yang sudah dipilih. 

\begin{exercisewrap}
\begin{nexercise} \label{raffleOf30TicketsWWOReplacement}
Departemen Anda mengadakan undian berhadiah. Mereka menjual 30 tiket dan menyediakan tujuh hadiah. (a) Tiket-tiket tersebut dimasukkan ke dalam topi dan satu tiket diundi untuk setiap hadiah. Tiket diambil tanpa pengembalian, yaitu tiket yang terpilih tidak dimasukkan kembali ke dalam topi. Berapakah peluang memenangkan hadiah jika Anda membeli satu tiket? (b)~Bagaimana jika tiket diambil dengan pengembalian?\footnotemark
\end{nexercise}
\end{exercisewrap}
\footnotetext{(a) Pertama, tentukan peluang tidak menang. Tiket diambil tanpa pengembalian, yang berarti peluang Anda tidak menang pada pengundian pertama adalah $29/30$, $28/29$ pada pengundian kedua, ..., dan $23/24$ pada pengundian ketujuh. Peluang Anda tidak memenangkan hadiah adalah hasil kali peluang-peluang tersebut: $23/30$. Artinya, peluang memenangkan hadiah adalah $1 - 23/30 = 7/30 = 0.233$. (b)~Ketika tiket diambil dengan pengembalian, terdapat tujuh pengundian yang saling independen. Sekali lagi, pertama-tama kita mencari peluang tidak memenangkan hadiah: $(29/30)^7 = 0.789$. Jadi, peluang memenangkan (setidaknya) satu hadiah ketika pengambilan dilakukan dengan pengembalian adalah 0.211.}

\begin{exercisewrap}
\begin{nexercise} \label{followUpToRaffleOf30TicketsWWOReplacement}
Bandingkan jawaban Anda dalam Latihan Terbimbing~\ref{raffleOf30TicketsWWOReplacement}. Seberapa besar pengaruh metode pengambilan sampel terhadap peluang Anda memenangkan hadiah?\footnotemark
\end{nexercise}
\end{exercisewrap}
\footnotetext{Peluang memenangkan hadiah sekitar 10\% lebih besar jika pengambilan sampel dilakukan tanpa pengembalian. Namun, dengan prosedur pengambilan sampel ini, paling banyak hanya satu hadiah yang dapat dimenangkan.}

Seandainya kita mengulangi Latihan Terbimbing~\ref{raffleOf30TicketsWWOReplacement} dengan 300 tiket, bukan 30 tiket, kita akan menemukan hal yang menarik: hasilnya akan hampir identik. Peluangnya adalah 0.0233 tanpa pengembalian dan 0.0231 dengan pengembalian. Ketika ukuran sampel hanya merupakan sebagian kecil dari populasi (di bawah 10\%), pengamatan hampir saling independen meskipun pengambilan sampel dilakukan tanpa pengembalian.


{\input{ch_probability/TeX/sampling_from_a_small_population.tex}}
